\chapter{LITERATURE REVIEW AND TECHNICAL BACKGROUND}

\section{Real-Time Tasks and Embedded Telemetry}
An RTOS partitions work into independently scheduled tasks and provides communication primitives that preserve ownership and ordering. FreeRTOS queues copy items into an internal buffer and are commonly used to transfer messages between producer and consumer tasks \cite{freertosQueues}. Mutexes add mutual exclusion and priority-inheritance behaviour for protected shared resources \cite{freertosMutex}. These mechanisms are directly relevant to the native KINETIQ control node: simulated sensing produces samples, classification consumes them, and telemetry publishes a single coherent output stream.

Real-time terminology requires restraint. A periodic task can measure release lateness and queue accumulation on a desktop simulator, but the host operating system and console I/O prevent that observation from constituting a hard timing guarantee. KINETIQ consequently records nominal period, observed lateness, queue high-water mark, dropped events, and missed-deadline count as diagnostics rather than advertising a real-time certification.

\section{Wearable Data and Android-Mediated Integration}
Wearable ecosystems often expose data through platform-controlled applications, accounts, permissions, and health-data APIs. Huawei documents Wear Engine as a development framework for selected wearable use cases \cite{huaweiWear}; its availability and access model do not establish an implemented direct integration in KINETIQ. The examined implementation instead relies on Huawei Health's visible user-interface data and a user-enabled Android accessibility service. Android's accessibility-service documentation establishes the controlled UI-observation context and permission model \cite{androidAccessibility}. It does not convert an accessibility capture into a medical-grade data interface.

This distinction leads to a useful evidence ladder. A value visible in Huawei Health is an upstream observation. A POST row at the local receiver establishes arrival at the receiver. A complete bridge output establishes eligibility for the eight-field mirror contract. A visible dashboard card establishes rendering. Each downstream claim requires its own evidence; an upstream screenshot cannot substitute for it.

\section{Validation, Freshness, and Provenance}
Network telemetry requires more than syntactically valid JSON. Field presence, range limits, timestamp interpretation, source identity, and receiver arrival time all shape whether an observation should control an interface. KINETIQ uses strict validation to reject missing or out-of-range values. Its UI then distinguishes an arriving value from an unchanged value and transitions through live, stale, and offline states. This is especially important for wellness presentation, where displaying a cached card without a status marker creates an avoidable ambiguity.

Data provenance is represented operationally rather than as an abstract label: the originating lane, simulated flag, transport class, timestamp, and link state determine whether a sample is accepted and how it is described. The current contract intentionally limits the shared mirror data to eight fields. Additional upstream observations, such as stress, remain documented at the receiver but cannot silently appear as mirror context unless the contract and forwarding logic are extended.

\section{Local Transport and Interface Security}
The Python standard library provides a basic HTTP server suitable for local development and controlled interfaces, while warning that production hardening remains the application's responsibility \cite{pythonHttp}. KINETIQ limits its examined dashboard server to \code{127.0.0.1} and uses an authenticated POST mechanism at the local receiver. This reduces accidental exposure on a local network but is not a complete security assessment. Secret storage, transport encryption across device boundaries, replay resistance, mobile application hardening, and independent penetration testing remain outside the documented work.

\section{Speech and Context-Grounded Assistance}
The voice pathway uses faster-whisper, an implementation associated with CTranslate2-based Whisper inference \cite{fasterWhisper}. A transcription latency on a synthesized clip measures only one constrained machine-and-input condition; it does not measure live microphone robustness, accent coverage, noise performance, or medical communication safety.

Grounding in KINETIQ is narrower than retrieval-augmented generation. Retrieval-augmented generation combines a parametric generator with external non-parametric memory \cite{lewis2020rag}; the examined KINETIQ assistant instead injects the accepted current telemetry snapshot into a bounded prompt or deterministic template. A local verifier checks whether key response elements agree with the accepted activity, exertion, heart-rate, steps, or PPG-quality context. The resulting verdict makes the implemented consistency check visible, but it cannot prove factual completeness or clinical validity.

\section{Embedded Hardware Context}
The ESP32 development-board documentation describes the board family and connection context used for the selected hardware \cite{espDevKit}. Espressif also documents the importance of bootstrapping pins and schematic precautions \cite{espHardware}. These sources support component-aware design but do not establish that KINETIQ completed physical assembly or acquired a final sensor stream. The repository includes a Wokwi simulation route; Wokwi documents its browser-based electronics simulation environment \cite{wokwi}. Simulation is useful for control logic and pin-planning checks, while physical sensor noise, power integrity, enclosure effects, and wearability require later validation.

\section{Engineering Gap}
Existing RTOS guidance explains task communication, and platform documentation explains constrained wearable or Android integration. Neither alone solves the presentation problem created when genuine observations, simulated control values, stale caches, and conversational output coexist in one mirror. KINETIQ addresses that bounded gap by making source and freshness explicit, applying a contract before UI and AI use, and retaining a labelled synthetic fallback. The work does not claim that no comparable product exists; it documents a transparent integration pattern for the specific implementation and evidence available.
